\centerline{\bf \Huge Целозначные многочлены}

\begin{flushright}
        \scriptsize{Эпизод 24. Последний посланник}
\end{flushright}

\textbf{Определение.}
Многочлен $P(x)$ с вещественным коэффициентами называется целозначным, если он принимает целые значения при всех целых $x$.

Рассмотрим многочлены вида

$$
\binom{x}{k}=C_x^k=\frac{x(x-1) \cdots(x-k+1)}{k!} .
$$


\problem{1}
Докажите, что многочлен $C_x^k$ --- целозначный.

\problem{2}
Дан многочлен $P(x)$ с вещественными коэффициентами степени $n$. Докажите, что $P(x)$ является линейной комбинацией многочленов $C_x^0, C_x^1, \ldots, C_x^n$ с вещественными коэффициентами, при этом коэффициенты определены однозначно.

\textbf{Определение.}
Пусть $P(x)$ --- многочлен с вещественными коэффициентами. \textit{Разностным многочленом} многочлена $P(x)$ называется многочлен

$$
\Delta P(x)=P(x+1)-P(x) .
$$


\problem{3}
Найдите многочлен $\Delta C_x^k$.

\problem{4}
Пусть $m$ --- целое число. Многочлен $P(x)$ с вещественными коэффициентами степени $n$ принимает целые значения в точках $m, m+1, \ldots, m+n$. Докажите, что этот многочлен является линейной комбинацией многочленов $C_x^0, C_x^1, \ldots, C_x^n$ с \textbf{целыми} коэффициентами. В частности, этот многочлен принимает целые значения во всех целых точках.
5. Многочлен $P(x)$ с вещественными коэффициентами степени $n$ таков, что числа $P\left(0^2\right), P\left(1^2\right)$, $P\left(2^2\right), \ldots, P\left(n^2\right)$ - целые. Докажите, что число $P\left(k^2\right)$ является целым для любого целого $k$.

\problem{6}
Даны простое число $p$ и целозначный многочлен $P(x)$. Для целого числа $n$ обозначим $r_n$ остаток от деления $P(n)$ на $p$. Докажите, что последовательность $\left\{r_n\right\}$ периодична.

\problem{7}
Целозначные многочлены $f(x)$ и $g(x)$ таковы, что для всех целых $n$ число $f(n)$ делится на $g(n)$. Докажите, что многочлен $f(x)$ делится на $g(x)$.